\paragraph{The Bag-Theoretic Structure of the Euclid--Mullin Dynamics}
\label{app:bag}
It is interesting to view the EM sequence as a dynamical system on bags of primes.
The state of the Euclid--Mullin construction at step~$n$ is fully captured
by the \emph{bag} (finite set) of primes collected so far:
$S_n = \{\seq{2}(0), \seq{2}(1), \ldots, \seq{2}(n)\}$.
The dynamics is a deterministic map on finite subsets of primes:
\[
S_{n+1} = S_n \cup \{\min \mathcal{F}(S_n)\},
\]
where $\mathcal{F}(S) = \text{PrimeFactors}(\prod S + 1)$ produces the
\emph{factor bag} of the Euclid number $E_n = \Prod{2}(n) + 1$.

The trajectory is entirely determined by $S_0 = \{2\}$.  The Markov
property holds: the future depends only on the current bag, not on the
order in which its elements were assembled.  The product
$\Prod{2}(n) = \prod S_n$ is a \emph{lossless encoding} of the bag (by
unique factorization of the squarefree integer~$\Prod{2}(n)$), and the
order of insertion is forgotten.

\paragraph{The pipeline: from bag to next prime.}
The selection of the next prime passes through a pipeline of operations,
each with distinct information-theoretic character:
\[
S_n \;\xrightarrow[\text{lossless}]{\prod(\cdot)}\;
\Prod{2}(n) \;\xrightarrow[+1]{\text{shift}}\;
\Prod{2}(n)+1 \;\xrightarrow[\text{factor}]{}\;
\mathcal{F}(S_n) \;\xrightarrow[\text{compress}]{\min}\;
\seq{2}(n+1).
\]

\emph{Stage~1: Product (lossless).}  The map $S_n \mapsto \Prod{2}(n)$ is
an injection (on sets of distinct primes), encoding ${\sim}\, n \log n$
bits of bag information into a single integer of ${\sim}\, 2^n$ digits.
This stage \emph{expands} the representation while preserving all
information.

\emph{Stage~2: The $+1$ shift (multiplicative $\to$ additive).}  The map
$\Prod{2}(n) \mapsto \Prod{2}(n)+1$ is the critical bridge between two worlds.
The integer~$\Prod{2}(n)$ lives in the multiplicative world: its structure
(as a product of known primes) is completely understood.  The
integer~$\Prod{2}(n)+1$ lives in the additive world: its factorization
bears no algebraic relationship to that of~$\Prod{2}(n)$, beyond the
guaranteed coprimality $\gcd(\Prod{2}(n),\, \Prod{2}(n)+1) = 1$.  This shift
is the \emph{source of mixing} in the dynamics.

\emph{Stage~3: Factorization (revealing the new bag).}  The factorization
$\Prod{2}(n)+1 = p_1^{a_1} \cdots p_k^{a_k}$ reveals the factor bag
$\mathcal{F}(S_n) = \{p_1, \ldots, p_k\}$.  This step is formally
lossless, but computationally opaque: predicting $\mathcal{F}(S_n)$ from
$S_n$ requires factoring a number of ${\sim}\, 2^n$ digits.

\emph{Stage~4: Minimum (catastrophic compression).}  The map
$\mathcal{F}(S_n) \mapsto \min \mathcal{F}(S_n)$ discards all but the
smallest element of the factor bag.  The compression ratio is
exponential: the input has ${\sim}\, 2^n$ bits; the output has $O(n)$
bits.  The composition of stages~2--4 is the \emph{minFac bottleneck}:
an information-destroying channel that maps ${\sim}\, 2^n$ bits to
$O(n)$ bits through arithmetically chaotic operations.

\paragraph{The orthogonal bag property.}
A fundamental structural fact distinguishes the factor bag
$\mathcal{F}(S_n)$ from the current bag~$S_n$:

\begin{proposition}[Orthogonal bags]
$\mathcal{F}(S_n) \cap S_n = \emptyset$.
That is, the primes dividing $\Prod{2}(n)+1$ are entirely disjoint from
the primes in the current bag.
\end{proposition}

\begin{proof}
For any $p \in S_n$: $p \mid \Prod{2}(n)$, hence
$\Prod{2}(n)+1 \equiv 1 \pmod{p}$, so $p \nmid \Prod{2}(n)+1$, i.e.\
$p \notin \mathcal{F}(S_n)$.
\end{proof}

This is the Euclidean coprimality at the heart of the construction, but
its bag-theoretic formulation reveals its dynamical significance: at each
step, the factorization produces a \emph{fresh sample} from the
complementary set $\mathbb{P} \setminus S_n$.  The new bag
$\mathcal{F}(S_n)$ contains no ``recycled'' primes---every element is
drawn from the universe of primes not yet seen.

Moreover, $\Prod{2}(n)$ is \emph{squarefree} (a product of distinct primes),
so $S_n$ is genuinely an unordered set with no multiplicities.  The
accumulator~$\Prod{2}(n)$ is the unique squarefree integer with prime
support exactly~$S_n$.

\paragraph{The CRT perspective.}
The Chinese Remainder Theorem provides the natural coordinate system for
the bag dynamics.  The integer $\Prod{2}(n)$ is determined by its CRT
vector:
\[
\mathbf{c}(n) = (\Prod{2}(n) \bmod 2,\;
\Prod{2}(n) \bmod 3,\; \Prod{2}(n) \bmod 5,\; \ldots)
\]
indexed by all primes.  For primes $p \in S_n$:
$\Prod{2}(n) \bmod p = 0$ (since $p \mid \Prod{2}(n)$).  For primes
$r \notin S_n$: $\Prod{2}(n) \bmod r = \prod_{p \in S_n} (p \bmod r)
\in (\ZZ/r\ZZ)^\times$.

The $+1$ shift maps $\mathbf{c}(n) \mapsto \mathbf{c}(n) + \mathbf{1}$
componentwise.  The factorization of $\Prod{2}(n)+1$ is determined by which
components satisfy $c_r(n) + 1 \equiv 0 \pmod{r}$, i.e., $c_r(n) = -1$.
The minFac is the \emph{winner of a race} across CRT coordinates: the
smallest prime $r \notin S_n$ for which $c_r(n) = -1$.  This race
depends on all coordinates simultaneously, but each coordinate is
determined independently (by CRT) from every other.

\paragraph{Fiber structure and non-reconstructibility.}
A natural question is whether the output of minFac carries information
about the bag that produced it.  The answer is: almost none.

\begin{proposition}[Massive many-to-one]
For any prime~$q$ and any set size~$k$, the number of bags $T$ of $k$
distinct primes satisfying $\minFac(\prod T + 1) = q$ grows
exponentially in~$k$.
\end{proposition}

The constraints on such a bag~$T$ are purely modular:
$\prod T \equiv -1 \pmod{q}$ (one congruence condition), plus
$\prod T \not\equiv -1 \pmod{p}$ for each prime $p < q$ with
$p \notin T$ (at most $\pi(q)$ non-congruence conditions).  The total
information content of the output is $O(q)$~bits, while the bag carries
${\sim}\, 2^k$~bits.  The fraction of information preserved is
$O(q / 2^k) \to 0$ super-exponentially.  Consequently, the map
$S \mapsto \minFac(\prod S + 1)$ is an \emph{exponentially lossy
summarization} of the bag.

\paragraph{MinFac as a greedy algorithm.}
The minimum operation in Stage~4 admits a natural optimization
interpretation.  Among all prime factors of $\Prod{2}(n)+1$, the minFac
rule selects the one that increases $\log \Prod{2}$ by the least amount:
\[
\seq{2}(n+1) = \arg\min_{p \in \mathcal{F}(S_n)} \log p.
\]
This is a \emph{greedy algorithm for slow growth}: at each step, the
accumulator grows as slowly as possible (given the constraint of
selecting a prime factor of $\Prod{2}(n)+1$).

Slow growth has a dynamical consequence.  Since $\Prod{2}(n)$ grows slowly
(relative to the maxFac alternative), each target prime~$q$ receives
\emph{more trials}---more steps during which $\Prod{2}(n)+1$ might be
divisible by~$q$---before the accumulator becomes so large that the
probability of $q$-divisibility becomes negligible.  The minFac rule
maximizes the number of opportunities each prime has to enter the
sequence.

By contrast, the maxFac rule (selecting the largest prime factor) causes
$\Prod{2}(n)$ to grow so rapidly that small primes quickly lose any chance
of being selected, which is why the second Euclid--Mullin sequence
provably misses infinitely many primes
(Cox--van der Poorten~\cite{CoxVdP1968}, Booker~\cite{BoSa2012}).

\paragraph{The $0$-accessibility of minFac.}
We formalize the advantage of minFac over other selection rules.  Define
a selection rule $\sigma\colon 2^{\mathbb{P}} \to \mathbb{P}$ mapping a
nonempty set of primes to one of its elements.  Say $\sigma$ is
\emph{$0$-accessible} for a prime~$q$ past the sieve gap if: whenever
$q \in F$ and all primes $p < q$ have been screened, $\sigma(F) = q$.

MinFac is $0$-accessible for every~$q$: past the sieve gap, all primes
below~$q$ are excluded from $\mathcal{F}(S_n)$, so if
$q \in \mathcal{F}(S_n)$, then $q = \min \mathcal{F}(S_n)$.  The event
``$q$ divides $\Prod{2}(n)+1$'' is \emph{sufficient} for~$q$ to enter the
sequence.

MaxFac is not $0$-accessible: even if $q \mid \Prod{2}(n)+1$, a larger
factor is selected instead, and the probability that~$q$ is the largest
factor of a number of size ${\sim}\, 2^{2^n}$ decays faster than any
polynomial.

The $0$-accessibility of minFac ensures that the walk mod~$q$ hitting
$-1$ is both necessary and sufficient for~$q$ to enter the sequence
(past the sieve gap).  This reduces the question of whether~$q$ appears
to a question about the walk alone, with no additional selection barrier.

\paragraph{Deterministic pseudorandomness.}
The EM sequence is fully deterministic, yet exhibits pseudorandom
behavior.  The pseudorandomness arises from the composition of two
sources of arithmetic opacity:
\begin{enumerate}[nosep]
\item \emph{The $+1$ shift} moves from the multiplicative world (where
  $\Prod{2}(n)$ is fully understood) to the additive world (where
  $\Prod{2}(n)+1$ is arithmetically opaque).  The relationship between the
  multiplicative structures of consecutive integers is the central
  mystery of analytic number theory.
\item \emph{The factoring step} reveals the prime structure of
  $\Prod{2}(n)+1$, but this structure is computationally hard to predict
  from~$\Prod{2}(n)$.  The computational difficulty of factoring is, in a
  precise sense, the \emph{source of pseudorandomness} for the EM
  sequence.
\end{enumerate}
Small changes to the bag (adding one prime~$p$) transform $\Prod{2}(n)$ to
$\Prod{2}(n) \cdot p$ and hence $\Prod{2}(n)+1$ to $\Prod{2}(n) \cdot p + 1$, a
completely different integer with a completely different factorization.
The dynamics is \emph{sensitive to the bag contents}, analogous to
sensitive dependence on initial conditions in chaotic systems, but
operating through number-theoretic rather than geometric mechanisms.

The orthogonal bag property reinforces this pseudorandomness: at each
step, the factor bag is drawn from $\mathbb{P} \setminus S_n$, a set
that has never been ``used'' in the construction.  Each step provides a
genuinely fresh arithmetic input, preventing the buildup of systematic
correlations.
